\problemname{Efterlyst}
\noindent
Police constable Acsel needs your help with an urgent matter, namely to catch the dangerous criminal Waxel.
Waxel is hiding somewhere in a city consisting of $N$ different locations, numbered from $1$ to $N$, with $M$ bidirectional roads each connecting two different locations.
He was for a long time at a certain location $X$ (it is unknown which), but then he moved to another location $Y$ (which we also do not know) by traveling along a sequence of roads.

The police have collected $K$ witness statements from people who saw Waxel.
Therefore they know that Waxel visited the locations $a_1, a_2, \dots, a_K$ on the way from $X$ to $Y$ (it is also possible that $a_i = X$ or $a_i = Y$ for some $i$).
However, they do not know the order in which the locations were visited.
Waxel may also have visited more than these $K$ locations on the way from $X$ to $Y$.

Your task is now to help the police find the locations that could possibly be $Y$, given that Waxel took a \textbf{shortest sequence of roads}\footnote{This means that the sum of the lengths of the roads that Waxel traveled along is as small as possible.} from $X$ to $Y$.
It is possible that there are no such locations at all, if the witness statements are not consistent with any shortest path.

\section*{Input}
The first line contains three integers:
\begin{itemize}
\item the number of locations in the city, $N$ ($2 \leq N \leq 2\cdot 10^5$),
\item the number of roads in the city, $M$ ($1 \leq M \leq 2 \cdot 10^5$), and
\item the number of witness statements, $K$ ($1 \leq K \leq N$).
\end{itemize}

The second line contains the $K$ \emph{distinct} integers $a_1, \dots, a_K$ ($1 \leq a_i \leq N$), the locations that Waxel visited.

The following $M$ lines describe the roads in the city.
The $i$th line contains three integers $u_i$, $v_i$ ($1 \leq u_i \not= v_i \le N$) and $w_i$ ($1 \leq w_i \leq 10^9$), meaning that the $i$th road connects locations $u_i$ and $v_i$ and has length $w_i$ meters.
It is guaranteed that it is possible to travel between any two locations by following a sequence of roads, and that there is at most one road connecting each pair of locations.

\section*{Output}
On the first line, print the number of nodes that could be $Y$.
Note that this number can be $0$.

On the second line, print all nodes that could be $Y$.
These should be printed separated by spaces in \textbf{increasing} order.

\section*{Scoring}
Your solution will be tested on a set of test groups.
To earn points for a group, you must pass all test cases in that group.

\noindent
\begin{tabular}{| l | l | p{10cm} |}
\hline
Group & Points & Constraints \\ \hline
$1$    & $5$         & There are exactly $N - 1$ roads, connecting locations $1$ and $2$, $2$ and $3$, and so on up to $N - 1$ and $N$ \\ \hline
$2$    & $16$         & There is only one sequence of (non-repeating) roads between each pair of locations, which implies that $M = N - 1$ \\ \hline
$3$    & $15$         & $N,M \leq 100$ \\ \hline
$4$    & $17$         & $N,M \leq 1\,000$ \\ \hline
$5$    & $15$         & $N,M \leq 10^5$ and $K \leq 5$ \\ \hline
$6$    & $32$         & No additional constraints \\ \hline
\end{tabular}

\section*{Explanation of Samples}
In sample $1$, the locations $1,2,4,5$ are possible destinations for Waxel.
To reach $2$, he could have taken the path $5-6-1-2$.

In sample $2$, there is no shortest path that passes through the given nodes.
The answer is thus $0$.

In sample $3$, there are quite many possibilities.
To reach $2$, Waxel could for example have taken the path $6-7-5-3-1-2$.
